% Generated draft: requires mathematical and source review.
\documentclass[11pt]{article}
\usepackage[margin=25mm]{geometry}
\usepackage{fontspec}
\setmainfont{Latin Modern Roman}
\usepackage{amsmath,amssymb,mathtools,unicode-math}
\setmathfont{Latin Modern Math}
\usepackage{enumitem,needspace,xurl,hyperref}
\setlength{\parindent}{0pt}
\setlength{\parskip}{5pt}
\setlength{\emergencystretch}{3em}
\newcommand{\N}{\mathbb{N}}
\newcommand{\Z}{\mathbb{Z}}
\newcommand{\Q}{\mathbb{Q}}
\newcommand{\R}{\mathbb{R}}
\newcommand{\C}{\mathbb{C}}
\newcommand{\F}{\mathbb{F}}
\newcommand{\A}{\mathbb{A}}
\newcommand{\PP}{\mathbb{P}}
\newcommand{\E}{\mathbb{E}}
\newcommand{\eps}{\varepsilon}
\newcommand{\dd}{\,\mathrm{d}}
\begin{document}
\section*{UnsolvedMath ID: 20000529}
\textbf{Source problem code:} catalogue entry
\subsection{Definitions and mathematical statement}
Fix \(d\geq1\) and a finite real exponent \(2<q<\infty\). For a Lebesgue measurable set \(E\subset\R^d\) of measure \(1\), use
\[
\widehat{\mathbf1_E}(\xi)=\int_E e^{-2\pi i x\cdot\xi}\,\dd x,
\qquad M_{d,q}=\sup_{|E|=1}\|\widehat{\mathbf1_E}\|_{L^q(\R^d)}.
\]
Christ's Question 2 asks to describe all measurable extremizers \(E\) satisfying \(\|\widehat{\mathbf1_E}\|_q=M_{d,q}\). Sets must be identified up to null sets; translations and volume-preserving invertible affine changes give the natural symmetries. An ellipsoid here is an invertible affine image of a ball, with volume \(1\). In one dimension this is an interval of length \(1\). The question is not a classification of arbitrary complex-valued functions under a norm constraint.

The source's Theorem 0.1, attributed to Christ [3], reports existence of an extremizing set for every \(q>2\); ellipsoid rigidity in each dimension for sufficiently large exponents sufficiently close to even integers; the corresponding rigidity in dimension one near even integers; and rigidity in dimension two for \(q\) sufficiently close to \(4\). The source writes the even exponents as \(2\N\) without giving thresholds, neighborhood sizes, or an explicit lower endpoint for that notation. Consequently no uniform neighborhood size in \(d\) or in the even exponent, and no extension to \(q=2\), is asserted here. Outside the supplied partial ranges, the global classification is described as open by the attached aid.

The aid supplies a one-dimensional test family
\[
E_g=[0,\tfrac12]\cup[\tfrac12+g,1+g],\qquad g\geq0.
\]
Writing \(t_+=\max(t,0)\) and
\(\delta(g)=\inf_{I:\,|I|=1}|E_g\mathbin{\triangle}I|\), it claims
\[
\delta(g)=\min(2g,1),\qquad
\|\widehat{\mathbf1_{E_g}}\|_4^4
=\frac12+\frac43(\tfrac12-g)_+^3,
\]
\[
\|\widehat{\mathbf1_{[0,1]}}\|_4^4-
\|\widehat{\mathbf1_{E_g}}\|_4^4
=\frac{\delta(g)}2-\frac{\delta(g)^2}2+\frac{\delta(g)^3}6.
\]
For each fixed finite \(q>2\), it also reports
\[
\lim_{g\to\infty}
\frac{\|\widehat{\mathbf1_{E_g}}\|_q^q}{\|\widehat{\mathbf1_{[0,1]}}\|_q^q}
=c_q=\frac{2\Gamma((q+1)/2)}{\sqrt\pi\,\Gamma((q+2)/2)}<1,
\]
and
\[
c_q=1-(\log2-\tfrac12)(q-2)+O((q-2)^2)
\quad(q\downarrow2).
\]
These supplied special-family claims compare equal half-intervals with a single interval. They do not classify all unit-measure sets, prove ellipsoid extremality for every \(q\), or give a global quantitative stability theorem.
\subsection{Short English statement}
For each dimension and finite \(q>2\), classify unit-measure sets maximizing the \(L^q\) norm of their indicator's Fourier transform. The source reports ellipsoid rigidity in specified neighborhoods of even exponents. The attached two-interval formulas are partial benchmarks rather than a global classification.
\subsection{Sources}
\begin{itemize}
\item Christ, Question 2 and Theorem 0.1, AIM workshop \emph{Carleson theorems and multilinear operators}: \url{https://aimath.org/pastworkshops/multilinopsproblems.pdf}. The theorem is attributed to Christ [3], but the full bibliography entry [3] is not supplied.
\item Local registry \texttt{data/unsolvedmath/problems.json}, permanent ID \(20000529\), exact source code \texttt{catalogue entry}; canonical \texttt{aim-analysis-notes.json}, notes[60].
\item The attached research attempt supplies the explicit gap and splitting formulas without a separate literature citation. Attempt 1 is explicitly machine-generated, not peer reviewed, and requires independent verification before citation.
\end{itemize}
\end{document}
